% ---------------------------------------------------------------------------
% Venue-neutral conference manuscript.
% To retarget: swap \documentclass for the venue template (e.g. aamas.cls,
% IEEEtran, or sigconf) -- the body text uses no class-specific commands.
% ---------------------------------------------------------------------------
\documentclass[10pt,twocolumn]{article}

\usepackage[margin=0.85in]{geometry}
\usepackage{amsmath,amssymb,amsthm}
\usepackage{graphicx}
\usepackage{booktabs}
\usepackage{microtype}
\usepackage[hidelinks]{hyperref}
\usepackage{caption}
\captionsetup{font=small,labelfont=bf}

\newtheorem{proposition}{Proposition}
\newtheorem{remark}{Remark}

\title{\bfseries A Selection-Conditioned Bias in the Evaluation of\\
Merit-Based Sortition, and What It Invalidates}

\author{%
  [Author Name]\\
  \small [Institution] \\
  \small \texttt{[email]}
}
\date{}

\begin{document}
\maketitle

\begin{abstract}
Merit-based sortition selects an active subset of participants in a decentralized
network by ranking an exponential moving average (EMA) of observed quality, with a
replacement percentile $P$ controlling how optimistically unobserved participants are
credited. We report that the quality metric used to evaluate this family of mechanisms
is biased: it scores the active set \emph{after} selection using the same observations
that drove the selection, so the recorded value is conditioned on the selection event.
The resulting bias is proportional to participant turnover and therefore to $P$ itself
--- the very parameter under study. We measure the bias at $0.108\times$ the per-epoch
churn rate ($R^2=0.874$), vanishing when churn vanishes.

The consequences are not cosmetic. Under the biased metric the apparent
quality-optimal percentile is $P^\star\in[70,85]\%$; under an unbiased metric it is
$P^\star\in[20,30]\%$, recovering the $[20,40]\%$ range reported by
Kruijssen et~al.~(2024) and reversing the ordering of competing mechanisms --- an
adaptive rule operating in $[55,85]\%$ ranks \emph{best} under the biased metric and
near-worst under the unbiased one. We then re-examine volatility-adaptive replacement
(VAMS), the natural extension of the framework. Across four pool dynamics, $100$ seeds
and both metrics, adaptive $P$ is statistically \emph{equivalent} to a fixed $P$ under
two one-sided tests. We give the mechanism: the control signal $\sigma_A$ varies at
exactly the sampling noise floor of a standard-deviation estimator
($0.0208$ measured against $0.0208$ predicted), so it carries no state information.
An oracle search over time-varying schedules for $P$, tuned with hindsight on the
data it is scored on, yields no gain over the best constant $P$ ($p=0.63$, $p=0.55$),
bounding the entire family of such policies.

Finally we show the framework is not at its ceiling: replacing the scalar percentile
with a per-participant uncertainty bonus improves quality over a correctly tuned static
baseline in all four pool dynamics, by up to $24\%$ of the merit-over-random gap
(paired $\Delta=+0.0118$, 95\% CI $[+0.0097,+0.0139]$).
\end{abstract}

% ===========================================================================
\section{Introduction}

Decentralized networks that reward measurable performance must repeatedly choose an
active subset of participants from a larger pool. Random sortition gives
representativeness but ignores quality. Kruijssen et~al.~\cite{kruijssen2024} propose
\emph{merit-based sortition}: participants are ranked by an EMA of their observed
quality, and a replacement percentile $P$ determines the value imputed to participants
who were not observed, which in turn governs how quickly they are promoted.

The framework invites an obvious extension. Real networks are non-stationary, so $P$
ought to adapt. We set out to evaluate exactly that hypothesis and found it false.
More importantly, we found that the instrument used to evaluate it --- the quality
metric --- is itself biased in a way that is correlated with the parameter under study,
which means the measurement cannot answer the question it is being asked.

This paper makes four contributions.

\begin{enumerate}\itemsep2pt
\item \textbf{A selection-conditioned bias} in the standard evaluation metric
(Section~\ref{sec:bias}), measured at $0.108\times$ churn, $R^2=0.874$.

\item \textbf{What it invalidates} (Section~\ref{sec:consequences}): the location of
the optimal $P$ moves from $[20,30]\%$ to $[70,85]\%$, and the ranking of competing
mechanisms inverts.

\item \textbf{Two negative results, stated as equivalences rather than failures to
reject} (Sections~\ref{sec:vams}--\ref{sec:headroom}): volatility-adaptive $P$ is
equivalent to fixed $P$ under TOST at $n=100$, with a measured mechanism for why; and
an \emph{oracle} time-varying schedule gains nothing, which bounds the whole policy
class rather than refuting one rule.

\item \textbf{A positive result} (Section~\ref{sec:ucb}): per-participant
uncertainty-directed selection beats a correctly tuned static baseline in all four
pool dynamics.
\end{enumerate}

We emphasise at the outset that the bias we report is a property of a
\emph{reimplementation} of the framework, and that measuring correctly
\emph{reproduces} the central conclusions of~\cite{kruijssen2024} rather than
overturning them. The cautionary finding is about how easily a plausible-looking
metric can invert a conclusion, not about the original mechanism.

% ===========================================================================
\section{Background}

\subsection{Merit-based sortition}
A pool of $N_{\mathrm{tot}}$ participants evolves over $N_{\mathrm{epochs}}$ epochs.
At each epoch $i$ an active set $A_i$ of size $N_{\mathrm{act}}$ is selected. Each
active participant $j$ produces an instantaneous quality
$T_{ij}\sim\mathcal{N}(\tilde Q_j, V)$, where the latent median quality
$\tilde Q_j\sim\mathcal{N}(\mu_Q,\sigma_Q)$ is fixed for the participant's lifetime.
Inactive participants produce nothing. A smoothed metric is maintained for every
participant,
\begin{equation}
Q_{ij} \;=\; \alpha\,\hat Q_{ij} + (1-\alpha)\,Q_{i-1,j},
\end{equation}
with $\hat Q_{ij}=T_{ij}$ for active participants and, crucially,
\begin{equation}\label{eq:impute}
\hat Q_{ij} \;=\; f_P\!\left(T_{i,A_i}\right)
\quad\text{for inactive } j,
\end{equation}
the $P$-th percentile of the active set's observed qualities. The next active set is
the top $N_{\mathrm{act}}$ participants by $Q_{ij}$.

\begin{remark}[$P$ is an exploration rate]\label{rem:explore}
Equation~\eqref{eq:impute} is the only place $P$ enters. It does not directly specify
how many participants are replaced; it sets how \emph{optimistically} unobserved
participants are credited. A high $P$ inflates the imputed value, lifting inactive
participants up the ranking and accelerating promotion. In bandit terms, $P$ is an
optimism bonus applied to unpulled arms --- a framing we return to in
Section~\ref{sec:ucb}.
\end{remark}

We use the parameters of~\cite{kruijssen2024} throughout: $\mu_Q=0.2$,
$\sigma_Q=0.1$, $V=0.2$, $\alpha=0.1$, $N_{\mathrm{epochs}}=1000$, and the four pool
scenarios in Table~\ref{tab:scenarios}.

\begin{table}[t]\centering\small
\caption{Pool scenarios, following~\cite{kruijssen2024}.}
\label{tab:scenarios}
\begin{tabular}{lrrrr}
\toprule
Scenario & $N_{\mathrm{init}}$ & $N_{\mathrm{act}}$ & $p_{\mathrm{growth}}$ & $p_{\mathrm{attr}}$\\
\midrule
Fixed     & 80  & 50 & $10^{-10}$ & $10^{-10}$\\
Growing   & 8   & 32 & $0.1$      & $10^{-10}$\\
Shrinking & 100 & 32 & $10^{-10}$ & $2\times10^{-3}$\\
Evolving  & 100 & 32 & $0.1$      & $2\times10^{-3}$\\
\bottomrule
\end{tabular}
\end{table}

\subsection{Experimental protocol}
All results use a vectorized engine validated against a direct transcription of the
reference implementation: across all four scenarios and both the fixed and adaptive
modes, the two engines agree (Welch $p>0.56$ in all eight comparisons, means within
$0.003$).

We report Welch's $t$-test, Hedges-corrected Cohen's $d$, and --- for every claim of
\emph{no difference} --- two one-sided tests (TOST) against a smallest effect size of
interest $\varepsilon=0.005$, which is $10\%$ of the merit-over-random quality gap.
This matters: at $n=10$ seeds, the protocol used in the reference notebooks, a
two-sample test has $80\%$ power only for $d\gtrsim1.3$. Absence of significance at
that sample size carries almost no information. We use $n=100$ throughout and $n=120$
paired seeds for the final comparison.

% ===========================================================================
\section{A selection-conditioned bias}
\label{sec:bias}

\subsection{Three estimands}
Let $\mathrm{lat}(S)=|S|^{-1}\sum_{j\in S}\tilde Q_j$ denote the mean latent quality of
a set $S$. The quantity a network operator cares about is the latent quality of the set
it fields. Three candidate measurements present themselves:
\begin{align}
M_1 &= \textstyle\frac{1}{|A_i|}\sum_{j\in A_i} T_{ij}, \\
M_2 &= \textstyle\frac{1}{|A_{i+1}|}\sum_{j\in A_{i+1}} \tilde T_{ij}, \\
M_3 &= \mathrm{lat}(A_{i+1}),
\end{align}
where in $M_2$, $\tilde T_{ij}=T_{ij}$ for retained participants
$j\in A_i\cap A_{i+1}$ and an independent fresh draw for newly promoted ones.
$M_2$ is the metric implemented in the reference notebooks; $M_3$ is ground truth,
available only in simulation.

\begin{proposition}\label{prop:bias}
$M_1$ is unbiased for $\mathrm{lat}(A_i)$, and in a stationary regime
$\mathbb{E}[M_1]=\mathbb{E}[M_3]$. $M_2$ is biased upward, with
$\mathbb{E}[M_2]-\mathbb{E}[M_3]=(1-c)\,b$, where $c$ is the churn rate and
$b=\mathbb{E}[T_{ij}\mid j\in A_{i+1}]-\tilde Q_j>0$.
\end{proposition}

\begin{proof}[Proof sketch]
Given $A_i$, each $T_{ij}$ is drawn independently of the selection rule, which has not
yet consumed it; hence $\mathbb{E}[M_1\mid A_i]=\mathrm{lat}(A_i)$. Stationarity of the
active-set composition gives $\mathbb{E}[\mathrm{lat}(A_i)]=\mathbb{E}[\mathrm{lat}(A_{i+1})]$.
For $M_2$, membership of $A_{i+1}$ is determined by $Q_{ij}$, which was just updated
using $T_{ij}$; conditioning on retention therefore selects on the realised
$T_{ij}$, so $\mathbb{E}[T_{ij}\mid j\in A_{i+1}]>\tilde Q_j$. Newly promoted
participants contribute independent draws and are unbiased, leaving the $(1-c)$ factor.
\end{proof}

The empirical check is decisive: $M_1$ and $M_3$ agree to four decimal places at every
$P$ in every scenario (Figure~\ref{fig:divergence}), exactly as
Proposition~\ref{prop:bias} requires, while $M_2$ departs from both.

\subsection{The bias scales with churn}
Proposition~\ref{prop:bias} leaves the magnitude of $b$ open, and $b$ is not constant:
raising $P$ inflates the imputed value in~\eqref{eq:impute}, intensifying competition
from the inactive pool, so retention demands a larger realised $T_{ij}$ and the
conditioning bias grows. Empirically the net effect is close to linear in churn and
passes through the origin (Figure~\ref{fig:churn}):
\begin{equation}
M_2-M_1 \;=\; 0.108\,c, \qquad R^2=0.874 .
\end{equation}
At $P\leq10\%$, where $c<0.003$, the two metrics are numerically identical; by
$P=100\%$ in the evolving pool, where $c=0.69$, they differ by $0.049$ --- comparable
to the entire merit-over-random effect the framework exists to produce.

\begin{figure}[t]\centering
\includegraphics[width=\columnwidth]{../paper_figs/fig1_metric_divergence.pdf}
\caption{Mean active-set quality against the replacement percentile. The unbiased
metric $M_1$ (blue) and latent ground truth $M_3$ (green, dashed) are
indistinguishable; the implemented metric $M_2$ (orange) diverges. Circles mark each
curve's argmax. Shading marks the $[20,40]\%$ range reported
by~\cite{kruijssen2024}, which contains the unbiased optimum in all four scenarios.}
\label{fig:divergence}
\end{figure}

\begin{figure}[t]\centering
\includegraphics[width=\columnwidth]{../paper_figs/fig2_bias_churn.pdf}
\caption{Left: the bias $M_2-M_1$ against per-epoch churn, pooled over all $P$ and all
four scenarios; the fit is through the origin. Right: churn is itself driven by $P$,
which is why the bias is correlated with the parameter under study.}
\label{fig:churn}
\end{figure}

% ===========================================================================
\section{What the bias invalidates}
\label{sec:consequences}

\subsection{The optimal percentile}
Under $M_2$ the apparent optimum is $P^\star=85,70,70,70\%$ for the fixed, growing,
shrinking and evolving pools. Under $M_1$ it is $P^\star=20,25,30,30\%$. The unbiased
optima fall inside the $[20,40]\%$ band reported in~\cite{kruijssen2024}, and the
unbiased curves reproduce the mechanism given there in words --- that at high
percentiles ``the active set is populated almost randomly'' --- which the biased
curves contradict.

\subsection{The ranking of mechanisms}
The reversal is not confined to tuning a scalar. Figure~\ref{fig:inversion} scores
eight mechanisms on both metrics in the evolving pool. Under $M_2$, an adaptive rule
operating over $[55,85]\%$ is the \emph{best} mechanism tested; under $M_1$ it is
second-worst, ahead only of random sortition ($d=-2.42$). A study using $M_2$ would
have selected and reported precisely the wrong mechanism.

\begin{figure}[t]\centering
\includegraphics[width=\columnwidth]{../paper_figs/fig3_ranking_inversion.pdf}
\caption{Eight mechanisms scored on both metrics (evolving pool, $n=100$, error bars
95\% CI). The metric determines which mechanism appears best.}
\label{fig:inversion}
\end{figure}

% ===========================================================================
\section{Volatility-adaptive $P$ is equivalent to fixed $P$}
\label{sec:vams}

\subsection{The hypothesis}
The natural extension of the framework makes $P$ respond to the network's current
volatility,
\begin{equation}\label{eq:vams}
P_t = P_{\min} + (P_{\max}-P_{\min})\,
      \varsigma\!\left(k\,(\sigma_{A,t-1}-\sigma_{\mathrm{ref}})\right),
\end{equation}
with $\varsigma$ the logistic function, $\sigma_{A,t}$ the standard deviation of the
active set's observed qualities and $\sigma_{\mathrm{ref}}$ its rolling mean over $W$
epochs. The intent is to replace more aggressively when the network is unstable.

\subsection{Result: equivalence, not merely non-significance}
Across four scenarios, $n=100$ seeds, and \emph{both} metrics,
equation~\eqref{eq:vams} is statistically equivalent to a fixed $P=25\%$ under TOST at
$\varepsilon=0.005$ (Table~\ref{tab:vams}). Re-centring the bounds on the per-scenario
optimum does not change this ($|d|\leq0.034$, TOST equivalent in all four). The null is
robust to the metric, to the adaptive range, and --- in a companion sweep --- to six
functional forms for~\eqref{eq:vams} and a $5\times5$ grid over $(k,W)$.

\begin{table}[t]\centering\small
\caption{Adaptive $P$ versus fixed $P=25\%$, $n=100$ seeds. All comparisons are
equivalent under TOST at $\varepsilon=0.005$.}
\label{tab:vams}
\begin{tabular}{lrrrr}
\toprule
& \multicolumn{2}{c}{on $M_1$} & \multicolumn{2}{c}{on $M_2$}\\
\cmidrule(lr){2-3}\cmidrule(lr){4-5}
Scenario & $d$ & $p$ & $d$ & $p$\\
\midrule
Fixed     & $-0.004$ & $0.975$ & $+0.010$ & $0.945$\\
Growing   & $-0.056$ & $0.691$ & $-0.035$ & $0.801$\\
Shrinking & $-0.009$ & $0.947$ & $+0.005$ & $0.974$\\
Evolving  & $-0.106$ & $0.453$ & $-0.077$ & $0.584$\\
\bottomrule
\end{tabular}
\end{table}

\subsection{Why: the control signal is noise}
The reason is measurable. For an active set of size $n$ with observation noise $V$, the
sampling standard deviation of a standard-deviation estimator is
$\approx\sigma/\sqrt{2n}$. In the fixed pool we measure the epoch-to-epoch standard
deviation of $\sigma_A$ at $0.0208$ against a predicted noise floor of $0.0208$; in the
evolving pool, $0.0265$ against $0.0258$. Because merit selection compresses the active
set's latent spread to $\approx0.057$ against $V=0.2$, $\sigma_A$ is dominated by
measurement noise and its fluctuations are statistically indistinguishable from pure
estimator noise.

Equation~\eqref{eq:vams} therefore steers $P_t$ with a noise process. The consequence
is visible in the realised control: $P_t$ has mean $24.99$ and standard deviation
$2.97$, lying within $\pm5$ points of $25\%$ in $91\%$ of epochs. VAMS is fixed-$P$
sortition plus jitter, by construction.

\begin{proposition}[Arity]\label{prop:arity}
No function of active-set summary statistics can supply the information the mechanism
needs. By Remark~\ref{rem:explore}, $P$ sets the optimism credited to unobserved
participants; the efficient optimism for participant $j$ depends on how much is known
about $j$ --- how long it has been unobserved. A scalar $P_t$ applies an identical
bonus to a participant idle for three epochs and one idle for three hundred. The
deficiency is in the dimension of the control variable, not the shape of its response.
\end{proposition}

Proposition~\ref{prop:arity} predicts the null before the experiment, and predicts that
no alternative functional form will rescue it --- which is what we observe.

% ===========================================================================
\section{A headroom bound on time-varying $P$}
\label{sec:headroom}

Refuting one rule at a time invites the objection that a better rule exists. We bound
the family instead. Consider two-phase schedules $P_t=p_1$ for $t<t_0$ and $p_2$
thereafter, and select the best $(p_1,p_2,t_0)$ \emph{using the same data it is scored
on}. This oracle has hindsight no causal rule can have, so its margin upper-bounds what
any online policy could achieve.

The oracle gains nothing: $+0.0013$ in the fixed pool ($d=0.096$, $p=0.63$) and
$-0.0016$ in the evolving pool ($d=-0.120$, $p=0.55$), both against the best constant
$P$. Within this framework, time-varying $P$ has no headroom to exploit.

% ===========================================================================
\section{Uncertainty-directed sortition}
\label{sec:ucb}

\subsection{Mechanism}
Proposition~\ref{prop:arity} points at the fix: make the optimism per-participant.
We retain the sortition framework and change two things. First, the estimator becomes
convergent --- a running sample mean over observed qualities rather than a
fixed-$\alpha$ EMA, which has a permanent error floor of
$V\sqrt{\alpha/(2-\alpha)}\approx0.046$ against a latent spread of only
$\sigma_Q=0.1$. Second, inactive participants receive no imputed observation at all;
instead selection ranks
\begin{equation}\label{eq:ucb}
s_j \;=\; \hat Q_j + c\sqrt{\frac{\log t}{n_j}},
\end{equation}
with $n_j$ the number of epochs $j$ has been active and $n_j=0$ given priority. This is
a combinatorial semi-bandit with $N_{\mathrm{act}}$ plays per round~\cite{chen2013}.

\subsection{Result}
Against a static baseline swept finely around its own optimum and compared on $120$
paired seeds using the unbiased metric, equation~\eqref{eq:ucb} wins in all four pool
dynamics (Table~\ref{tab:ucb}, Figure~\ref{fig:ucb}).

\begin{table}[t]\centering\small
\caption{Uncertainty-directed selection versus a correctly tuned static $P$, scored on
$M_1$, $120$ paired seeds.}
\label{tab:ucb}
\begin{tabular}{lrrrr}
\toprule
Scenario & best $P$ & $\Delta$ & 95\% CI & $d$\\
\midrule
Fixed     & $20\%$ & $+0.0040$ & $[0.0037,0.0042]$ & $0.34$\\
Growing   & $25\%$ & $+0.0057$ & $[0.0041,0.0073]$ & $0.37$\\
Shrinking & $25\%$ & $+0.0037$ & $[0.0034,0.0040]$ & $0.23$\\
Evolving  & $30\%$ & $+0.0118$ & $[0.0097,0.0139]$ & $0.93$\\
\bottomrule
\end{tabular}
\end{table}

\begin{figure}[t]\centering
\includegraphics[width=\columnwidth]{../paper_figs/fig4_ucb_result.pdf}
\caption{Left: paired improvement over a correctly tuned static $P$, 95\% CI; the
dashed line is the pre-registered smallest effect of interest. Right: the exploration
constant has a broad interior optimum.}
\label{fig:ucb}
\end{figure}

We report this carefully. All four intervals exclude zero, but in the fixed and
shrinking pools the gain is \emph{below} our pre-registered $\varepsilon=0.005$ and the
same data support an equivalence verdict; the improvement there is reliable but
practically marginal. Scaled against the merit-over-random gap of $\approx0.05$, the
gains are $8\%$, $11\%$, $7\%$ and $24\%$. Only in the evolving pool --- the most
non-stationary scenario, where per-participant uncertainty varies most --- is the
effect substantial. The exploration constant is not knife-edged but it does matter:
$c\in[0.15,0.3]$ improves on the tuned static baseline in every scenario, with an
interior optimum near $c\approx0.2$, whereas $c\leq0.05$ under-explores and loses to it
in half the scenarios ($-0.0060$ fixed, $-0.0006$ shrinking at $c=0.05$).
Equation~\eqref{eq:ucb} therefore trades one tuned scalar ($P$) for another ($c$);
its advantage is that $c$ modulates a per-participant quantity, which
Proposition~\ref{prop:arity} identifies as the information $P$ cannot carry.

\subsection{Estimator memory}
Sweeping the estimator's effective smoothing against a mean-reverting drift in latent
quality (with stationary variance held fixed, so drift is not confounded with a
widening quality distribution) shows the reference value $\alpha=0.1$ is never optimal:
the best effective smoothing is $\approx0.02$, a memory of roughly $50$ epochs rather
than $10$, at every drift rate tested. A small amount of forgetting does beat infinite
memory once drift is present, so the trade-off is real --- it is simply mis-calibrated
by about a factor of five.

% ===========================================================================
\section{Discussion}

\paragraph{The general lesson.}
The bias reported here is not exotic. It arises from a natural-looking choice --- score
the set you just selected, using the observations you have --- whose defect is that the
selection consumed those observations. The dangerous feature is not its size but its
\emph{correlation with the independent variable}: because churn rises with $P$, the
bias masquerades as a treatment effect. Any mechanism that increases turnover is
flattered. We suggest a simple discipline for this class of evaluation: score the
population \emph{before} the selection rule consumes the observations, and verify
against latent ground truth where simulation permits.

\paragraph{Limitations.}
First, our results concern the simulation model of~\cite{kruijssen2024}, in which
latent quality is fixed for a participant's lifetime; conclusions about estimator
memory in particular may not transfer to networks with substantial drift, though our
OU-drift sweep probes this directly. Second, Proposition~\ref{prop:bias} is stated for
a stationary regime and verified numerically rather than proved in full generality.
Third, we adapt the regret framework of~\cite{chen2013} by analogy and do not derive a
bound specific to this mechanism; the sortition setting differs in that the
\emph{pool itself} grows and shrinks. Fourth, all results are simulation-only;
validation on deployed network data remains open.

\paragraph{Future work.}
With a correct measurement harness in place, the search for mechanisms that genuinely
improve on tuned static sortition is well posed for the first time. The gains we report
are modest, and Section~\ref{sec:headroom} suggests the remaining headroom in
\emph{scalar} policies is near zero --- so progress must come from richer control, of
which~\eqref{eq:ucb} is the simplest instance.

% ===========================================================================
\section{Conclusion}

The metric used to evaluate merit-based sortition conditions on the selection event it
is meant to evaluate, producing an upward bias proportional to turnover and therefore to
the replacement percentile itself. Correcting it moves the apparent optimal percentile
from $[70,85]\%$ to $[20,30]\%$ --- reproducing, rather than overturning, the range
reported by Kruijssen et~al. --- and reverses which mechanisms appear best. Measured
correctly, volatility-adaptive replacement is equivalent to fixed replacement, for a
reason we can state and measure: its control signal sits at the noise floor, and a
scalar control cannot carry per-participant information in any case. An oracle schedule
confirms the whole family is exhausted. Replacing the scalar with a per-participant
uncertainty bonus does improve on a correctly tuned baseline, substantially so in the
most dynamic regime.

% ===========================================================================
\begin{thebibliography}{9}\small

\bibitem{kruijssen2024}
J.~M.~D. Kruijssen, R.~Valieva, K.~Peluso, N.~Emmons, and S.~N. Longmore.
\newblock Merit-Based Sortition in Decentralized Systems.
\newblock \emph{Allora Decentralized Intelligence}, 1:20--27, 2024.
\newblock arXiv:2411.07302.

\bibitem{chen2013}
W.~Chen, Y.~Wang, and Y.~Yuan.
\newblock Combinatorial multi-armed bandit: General framework and applications.
\newblock In \emph{Proc. ICML}, 2013.

\bibitem{auer2002}
P.~Auer, N.~Cesa-Bianchi, and P.~Fischer.
\newblock Finite-time analysis of the multiarmed bandit problem.
\newblock \emph{Machine Learning}, 47:235--256, 2002.

\bibitem{garivier2011}
A.~Garivier and E.~Moulines.
\newblock On upper-confidence bound policies for switching bandit problems.
\newblock In \emph{Proc. ALT}, 2011.

\bibitem{lakens2017}
D.~Lakens.
\newblock Equivalence tests: A practical primer for $t$-tests, correlations, and
meta-analyses.
\newblock \emph{Social Psychological and Personality Science}, 8(4):355--362, 2017.

\bibitem{flanigan2021}
B.~Flanigan, P.~G\"olz, A.~Gupta, B.~Hennig, and A.~D. Procaccia.
\newblock Fair algorithms for selecting citizens' assemblies.
\newblock \emph{Nature}, 596:548--552, 2021.

\end{thebibliography}

\end{document}
